\subsection{W-7 (BIG)}
Principy výpočtů v distribuovaných systémech - Spark framework. Principy distribuované indexace - Elastic.

\subsubsection*{Apache Spark – principy distribuovaného výpočtu}

\textbf{Apache Spark} je open-source výpočetní framework pro distribuované zpracování dat. Hlavní výhody oproti MapReduce (Hadoop) zahrnují rychlost (díky in-memory zpracování) a jednoduché API pro práci s daty.

\paragraph{Architektura Spark}
Spark pracuje na principu \textbf{Master-Worker} architektury:
\begin{itemize}
    \item \textbf{Driver Program} – řídí aplikaci, plánuje úlohy, spravuje metadata.
    \item \textbf{Cluster Manager} – přiděluje prostředky (např. YARN, Mesos, Kubernetes).
    \item \textbf{Executors} – vykonávají úlohy, zpracovávají data.
\end{itemize}

\paragraph{RDD (Resilient Distributed Dataset)}
Základní datová abstrakce ve Sparku. RDD je:
\begin{itemize}
    \item \textbf{Imutabilní} – neměnné.
    \item \textbf{Paralelně distribuované} – rozdělitelné mezi uzly.
    \item \textbf{Odolné vůči chybám} – díky \textbf{lineage} (zachování historie operací).
\end{itemize}

\paragraph{Transformace vs Akce}
\begin{itemize}
    \item \textbf{Transformace} (např. \texttt{map}, \texttt{filter}) – definují nové RDD, ale nespouští výpočet (jsou \textbf{lazy}).
    \item \textbf{Akce} (např. \texttt{collect}, \texttt{count}) – spustí výpočet a vrací výsledek.
\end{itemize}

\paragraph{DAG – Directed Acyclic Graph}
Spark interně reprezentuje výpočetní úlohu jako DAG – graf, kde uzly představují RDD a hrany transformace. Tento DAG je optimalizován a rozdělen na etapy (stages).

\paragraph{Shuffle operace}
Při operacích jako \texttt{groupByKey} nebo \texttt{reduceByKey} dochází ke \textbf{shuffle} – přerozdělení dat mezi uzly, což je náročné na čas i zdroje.

\paragraph{Odolnost vůči chybám}
Spark nepoužívá tradiční replikaci jako Hadoop, ale obnovuje data pomocí \textbf{lineage}, tedy znovuvýpočtem původních transformací.

\paragraph{Příklad jednoduchého programu}
\begin{verbatim}
val data = sc.parallelize(Seq(1, 2, 3, 4))
val result = data.map(x => x * 2).reduce(_ + _)
\end{verbatim}

\subsubsection*{Elasticsearch – principy distribuované indexace}

\textbf{Elasticsearch} je vyhledávací engine postavený nad knihovnou \textbf{Lucene}. Umožňuje distribuované indexování a vyhledávání dokumentů.

\paragraph{Základní pojmy}
\begin{itemize}
    \item \textbf{Cluster} – skupina uzlů, které spolupracují.
    \item \textbf{Node} – jeden server v clusteru.
    \item \textbf{Index} – kolekce dokumentů, podobné databázi.
    \item \textbf{Document} – základní jednotka dat (v JSON formátu).
    \item \textbf{Shard} – část indexu, umožňuje paralelizaci.
    \item \textbf{Replica} – kopie shardu pro zvýšení dostupnosti a výkonu.
\end{itemize}

\paragraph{Inverzní index}
Základní struktura pro vyhledávání v textu:
\begin{itemize}
    \item Umožňuje rychlé vyhledání dokumentů obsahujících určitý termín.
    \item Místo uložení dokumentů jako celků uchovává mapování termín → seznam dokumentů.
\end{itemize}

\paragraph{Distribuované indexování}
Při vložení dokumentu:
\begin{itemize}
    \item Elasticsearch podle hash funkce určí primární shard.
    \item Dokument je indexován a distribuován mezi uzly.
    \item Repliky jsou vytvořeny podle nastavení (např. 1 replika = 2 kopie dokumentu).
\end{itemize}

\paragraph{Dotazovací jazyk (Query DSL)}
Elasticsearch poskytuje flexibilní dotazovací jazyk v JSON formátu:
\begin{verbatim}
GET /index/_search
{
  "query": {
    "match": {
      "title": "Distribuované systémy"
    }
  }
}
\end{verbatim}

\paragraph{Vyhledávání ve více šardech}
Dotaz je paralelně poslán na všechny primární a replikované shardy. Výsledky jsou agregovány a seřazeny podle relevance.

\subsubsection*{Shrnutí}
Spark a Elasticsearch představují dva zásadní přístupy k distribuovanému zpracování – Spark se zaměřuje na výpočty, zatímco Elasticsearch na indexaci a vyhledávání. Oba systémy využívají shardování, paralelismus, a odolnost vůči chybám jako základní stavební prvky distribuovaných architektur.

